\documentclass[journal,twocolumn]{IEEEtran}
% Working draft 2026-09-11: see IEEE_PAPER_HANDOFF.md for evidence and open checks.
% Existing main.pdf is stale. Verify author details, dataset rights and venue format.

\usepackage{cite}
\usepackage{amsmath,amssymb,amsfonts}
\usepackage{graphicx}
\usepackage{booktabs}
\usepackage{hyperref}

\begin{document}

\title{AITE: A Reproducible Prototype for Isolated ASL Recognition and Recorded ISL Pose Playback}

\author{Monisha Sharma, Nandita R Nadig, Naman Nagar, and Dr. Sudhamani M J}

\maketitle

\begin{abstract}
This paper describes AITE, a local research prototype that accepts a short,
isolated American Sign Language (ASL) video, estimates body and hand landmarks,
classifies one ASL gloss, and presents recorded two-dimensional Indian Sign
Language (ISL) pose motion where a configured recording lookup is available.
The recognition component is a temporal Transformer trained and evaluated with
signer-disjoint splits of an MSASL subset. The selected 100-class model achieved
47.02\% Top-1 accuracy, 76.78\% Top-5 accuracy, and 44.61\% macro-F1 on a
held-out test split. A 126-class experimental expansion achieved 50.42\% Top-1
and 45.99\% macro-F1, but is not deployed because its additional ASL--ISL
playback mappings have not undergone ISL linguistic review. The system provides
historical engineering checks for three selected examples, not linguistic
validation or current browser acceptance. It is not a continuous
sign-language recognizer, a validated ASL--ISL machine-translation system, or a
pose-generation model. These boundaries are reported explicitly to make the
prototype reproducible and to prevent recorded retrieval from being represented
as generative signing.
\end{abstract}

\begin{IEEEkeywords}
Sign language recognition, ASL, ISL, pose landmarks, Transformer, reproducibility
\end{IEEEkeywords}

\section{Introduction}
ASL and ISL are distinct natural languages; an English label match alone does
not establish equivalent lexical meaning, grammar, or motion. Consequently, a
cross-lingual sign system must evaluate recognition, translation, and signing
output separately. AITE is designed as a transparent prototype rather than an
end-to-end solution: it validates an isolated-sign recognition path and replays
available ISL landmark recordings with explicit provenance.

The contribution is practical and bounded: (i) a signer-disjoint landmark-based
isolated ASL recognition experiment, (ii) a local API and browser interface with
confidence rejection, and (iii) an ISL recorded-pose playback path that labels
unsupported outputs as illustrative rather than linguistic output. This work does
not claim sentence recognition, translation adequacy, or avatar intelligibility.

\section{System}
\subsection{Input and recognition}
The extraction design uses MediaPipe hand tracking~\cite{zhang2020hands} and
body landmarks sampled into 32 frames with 33
pose landmarks and 21 landmarks for each hand. Coordinates are normalized using
the shoulder midpoint and shoulder width. A temporal Transformer consumes 75
landmarks per frame along with velocity and landmark-presence features. It is
SPOTER-inspired~\cite{bohacek2022spoter}, not an exact reproduction. The API
returns a predicted gloss only when its softmax confidence exceeds a threshold
selected on the validation split; otherwise it returns an explicit low-confidence
rejection. The deployed model scope is one MSASL-100-style isolated sign per
clip, not a phrase or sentence.

\subsection{Draft gloss conversion and output}
The application defaults to deterministic, rule-based draft gloss conversion.
Optional Gemini and Llama refinement are disabled by default; no translation
accuracy gain is established. Token checks do not guarantee meaning preservation. It is not
a trained ASL-to-ISL translation model, and any gloss displayed in the user
interface requires signer review. The renderer looks up recorded ISL landmark
motion when a configured lookup matches. Otherwise it may produce a visibly
labelled illustrative two-dimensional animation. Thus, playback is retrieval of
recorded motion rather than pose generation or a realistic avatar. Multiword
fallback can select only one sign. Fingerspelling awaits real alphabet assets
and complete visual acceptance.

\section{Data and protocol}
The primary recognition experiment uses an MSASL~\cite{joze2019msasl} subset with signer-disjoint
training, validation, and test partitions. Landmark-quality filtering retained
2,747 training, 643 validation, and 521 test clips for the selected 100-class
model. The INCLUDE-50 ISL word collection supplies 61 recorded word-pose classes
for playback, while ISL-CSLRT supplies recorded sentence-pose sequences. Neither
collection provides verified parallel ASL--ISL sentence supervision.

The recorded protocol uses validation data for model selection and confidence
calibration, reserving the test partition for subsequent evaluation.
Reported selective accuracy is accuracy among predictions accepted above the
validation-selected threshold; coverage is the fraction of clips accepted. This
is important because rejection improves reliability only by declining some
inputs.

\section{Results}
\begin{table}[t]
\caption{Signer-disjoint isolated ASL recognition results. The 126-class result
is experimental and is not connected to the API.}
\label{tab:recognition}
\centering
\begin{tabular}{lccc}
\toprule
Model & Top-1 & Top-5 & Macro-F1 \\
\midrule
MSASL-100, validation & 49.30\% & 75.58\% & 45.34\% \\
MSASL-100, test & 47.02\% & 76.78\% & 44.61\% \\
MSASL-126 candidate, validation & 45.13\% & 72.43\% & 42.73\% \\
MSASL-126 candidate, test & 50.42\% & 76.39\% & 45.99\% \\
\bottomrule
\end{tabular}
\end{table}

The selected MSASL-100 threshold was 0.590399. On the held-out test partition,
accepted predictions reached 81.73\% selective accuracy at 37.81\% coverage.
The experimental 126-class model selected threshold 0.636551 on validation; it
reached 84.52\% selective accuracy at 28.33\% test coverage. Its expanded class
set should not be interpreted as expanded ASL-to-ISL support: linguistic
equivalence and ISL recorded-motion suitability remain unreviewed.

The saved helper-script report records three engineering examples using local ASL
videos: \textit{hello} to recorded \textit{NAMASKAR}, \textit{drink} to recorded
\textit{drink}, and \textit{man} to recorded \textit{man}. All other recognised
glosses require separate output validation. This historical report is not a new
current-API or browser acceptance test. The verification script, input paths, confidence, feature coverage,
and output mode are retained in the repository documentation.

\section{Limitations and ethics}
The two ASL experiments use different test populations, not a controlled
improvement comparison. An offline motion autoencoder obtained test MAE 0.40341
in normalized coordinates; it reconstructs poses and is not text-conditioned.
Code review identified an extraction-runtime mismatch between upload and the
historical helper, browser MIME compatibility risks, incomplete multiword output
and unvalidated approximate sentence lookup. The suite passed 38 tests, but
physical-camera and alphabet-asset acceptance remain pending.
The recognition metrics are moderate and apply only to the stated isolated-sign
data distribution. They do not measure continuous signing, webcam use in the
wild, sentence accuracy, ASL-to-ISL translation quality, non-manual markers, or
human comprehension of the output. Rule-based text reordering cannot establish
ISL grammar. A two-dimensional landmark playback is neither an ISL avatar nor a
pose-generation model. The prototype must not be used for safety-critical,
medical, legal, or high-stakes communication.

Before public or clinical use, the project requires licensed parallel
ASL--ISL data, ISL linguist/signer review of labels and output, evaluation by
Deaf users, continuous-sign recognition benchmarks, privacy and retention rules
for uploaded video, and a validated motion-generation or reviewed retrieval
library. The application should disclose confidence, provenance, and fallback
status for every result.

\section{Conclusion}
AITE demonstrates a reproducible local pipeline for signer-disjoint isolated
ASL recognition and limited recorded ISL pose playback. Its measured recognition
results and explicit rejection mechanism make it suitable for research demos
within this narrow scope. Completing an ASL-to-ISL translation product remains
dependent on reviewed cross-lingual supervision and evaluation of intelligible
ISL output; those requirements are not satisfied by the present prototype.

\bibliographystyle{IEEEtran}
\bibliography{paper_sources}

\end{document}
